\section{USB command and module interfaces}
There is no network protocol. USBX CDC ACM is a byte stream to a computer; USB packets do not frame commands. Proposed requests are ASCII lines, at most 128 bytes including LF, with optional CR immediately before LF. One request is outstanding per connection. Discard malformed or overlength input through the next LF; never execute a partial line. Parse unsigned values with explicit overflow/range checks.
\begin{SDDTable}{lY}
Request & Contract\\\midrule
\texttt{STATUS} & Versions, operating state, active ID, storage usage, durable count and per-path fault/drop counters.\\
\texttt{LIST} & Closed session IDs, generations, committed logical lengths and incomplete flags.\\
\texttt{START} / \texttt{STOP} & Allocate and begin a new session / complete the cutoff-and-drain handshake.\\
\texttt{READ id offset count} & Read 1--512 bytes from a closed session after validating its logical bounds and export lease.\\
\texttt{ERASE\_PREPARE id} & Validate a closed, unpinned session and issue a connection-bound confirmation token.\\
\texttt{ERASE\_CONFIRM id generation token} & Consume a matching token once; journal deletion intent and start bounded asynchronous erasure in maintenance mode.\\
\texttt{OP\_STATUS operation\_id} & Return deletion progress, verified completion, or failure.\\
\end{SDDTable}
Replies are LF-terminated \texttt{OK ...} or \texttt{ERR code ...}. A read returns \texttt{DATA id offset length crc32} plus LF, exactly \texttt{length} binary bytes, then LF. The chunk CRC uses the algorithm in Section~\ref{sec:persistence}. The host respects the length even when binary bytes contain LF. Reject invalid state, unknown ID, stale generation, bounds violations and busy/full resources with distinct error codes.

READ pins the session for a proposed 30 s inactivity lease. Release the pin on explicit connection close, disconnect or lease expiry after in-flight access is quiescent. A proposed 30 s erase token binds connection, session and generation; invalidate it on disconnect, expiry or use. Pinning and token validation occur in data so deletion cannot race export. Duplicate confirmation returns existing operation status while that operation is retained, rather than erasing twice.

The host client performs STATUS, STOP, LIST, repeated READ with CRC checks, file/manifest save, then a separate explicit deletion confirmation. Reconnect resumes only after rechecking session ID and generation. The build helper is a firmware programmer, not this client.

\subsection{USB adapter implementation and print capture}
The USB branch implements bounded parsing and queued command/response delivery; storage semantics remain owned by the data worker. Until that worker registers readiness, STATUS reports NOT\_READY and session operations return an explicit NOT\_READY error. The adapter never fabricates persisted sessions or calls flash SPI directly. Concrete text replies, queue ownership and integration rules are recorded in \texttt{docs/usb-integration.md}. LIST uses zero or more \texttt{OK SESSION id generation length incomplete} rows followed by \texttt{OK END}; generation and token arguments are unsigned decimal integers. STATUS includes a stable device UID for host resume checks.

Two transport extensions support diagnostic output without corrupting binary exports. \texttt{LOG\_READ cursor count} snapshots up to 512 bytes from an 8192-byte RAM print history. It replies with \texttt{LOG start length end overwritten crc32}, LF, exactly length bytes, then LF; a zero-length snapshot is permitted. Cursor positions wrap modulo \(2^{32}\), and expired cursors clamp to the oldest retained byte. The host saves original print bytes and capture metadata for later viewing. Board reset clears this RAM history; diagnostic prints are not persisted in sensor flash sessions. Standard output/error and the character-output hook append without waiting for USB; history overflow overwrites the oldest bytes and is counted.

\texttt{CLOSE} replies \texttt{OK CLOSE} and invalidates connection-bound work. USB deactivation or DTR deassertion also changes the connection generation. The data worker must observe this even without a queued command, then release old leases/tokens after in-flight access is quiescent. Transport failures and partial frames require close/reopen; they do not imply that an accepted storage operation was rolled back. The host CLI and optional Tk application are in \texttt{client/}; export verifies chunk CRCs, saves a manifest, and rechecks device/session identity plus saved chunks before resuming.
